\documentclass[12pt]{article}
\usepackage[utf8]{inputenc}
\usepackage[T1]{fontenc}
\usepackage{hyperref}
\title{Tsishingini Language Resources for K-12 Curriculum Development}
\author{Compiled by Claude Sonnet 5 (claude-sonnet-5) for Andr\'as Kornai}
\date{}
\begin{document}
\maketitle
\section{Introduction}
Tsishingini (ISO 639-3 tsw; also Salka, Shingini) is a Kainji (Kambari) language of Niger State, Nigeria, with about 100,000 speakers. Scripture and audio materials exist; no educational materials were found.
\section{Primary Educational Resources}
None found.
\section{Linguistic and Reference Materials}
A short overview is at \url{https://en.wikipedia.org/wiki/Shingini_language}. No dictionary or grammar was found.
\section{Religious and Community Resources}
A Tsishingini New Testament (2011, Wycliffe Bible Translators; facilitated by the Kambari Language Project) is at \url{https://dbs.org/bibles/TSWWBT}. Portions of the Old Testament (2025) are on eBible: \url{https://ebible.org/details.php?id=tsw}, \url{https://ebible.org/study/content/texts/tsw/about.html}. Scripture Earth lists further scripture, Jesus Film and Luke video: \url{https://www.scriptureearth.org/00eng.php?iso=tsw}.
\section{Digital and Audio Resources}
Global Recordings Network has audio and video in the language: \url{https://globalrecordings.net/en/language/tsw}.
\section{Curriculum Development Context}
The scripture texts are the main written corpus and show an established orthography, which could support early literacy work. Licences differ: the eBible Old Testament is CC BY-SA 4.0, while the New Testament is non-commercial and no-derivatives. No school programmes were found.
\end{document}